% Source: Solutions/Srijit-PS2-Q12-15.pdf, pp. 1-2 (typed). Restated problem statement omitted.
% REVIEW: (1) the proof tacitly assumes x1 < x2 (never stated; WLOG is needed).
%         (2) (-inf, x2) = U_{a<x2} (a, x2) fails when X has a least element m: m lies in the
%             ray but in no (a, x2). The rays are open anyway (they are subbasis elements of
%             the order topology), so the conclusion stands. Kept as written.
\submittedby{Srijit Bhadra}

\emph{Solution:} Let $X\neq\emptyset$ and $\mathscr{T}$ be the order topology on $X$. Let
$x_1,x_2\in X$ such that $x_1\neq x_2$. Consider the basic open set
$(x_1,x_2)=\{x\in X : x_1<x<x_2\}\in\mathscr{T}$. Now, there can be two cases:

\textbf{Case 1.} \underline{$(x_1,x_2)=\emptyset$.} Observe that
\[ (-\infty,x_2)=\bigcup_{a\in X,\ a<x_2}(a,x_2)\in\mathscr{T}. \]
Similarly,
\[ (x_1,\infty)=\bigcup_{b\in X,\ b>x_1}(x_1,b)\in\mathscr{T}. \]
Also, $x_1\in(-\infty,x_2)$ and $x_2\in(x_1,\infty)$. Crucially,
$(-\infty,x_2)\cap(x_1,\infty)=\emptyset$.

\textbf{Case 2.} \underline{$(x_1,x_2)\neq\emptyset$.} Take any point $c\in(x_1,x_2)$. As we saw
earlier, $(-\infty,c)$ and $(c,\infty)$ are open in $X$. Note that $x_1\in(-\infty,c)$,
$x_2\in(c,\infty)$, and $(-\infty,c)\cap(c,\infty)=\emptyset$.

Thus, in both cases, we found two disjoint open sets in $X$ with one containing $x_1$ and the
other containing $x_2$. As $x_1,x_2\in X$ were chosen arbitrarily, we can say that
$(X,\mathscr{T})$ is Hausdorff. Since $(X,\mathscr{T})$ was also chosen arbitrarily, we can say
that every order topology is Hausdorff. \hfill$\square$
